
\documentclass[a4paper]{article}
\usepackage[pdftex]{graphicx}
\usepackage[utf8]{inputenc}
\usepackage{enumerate}
\usepackage{icomma}
\usepackage{amssymb}
\usepackage{amsmath}
\usepackage{tikz}
\usepackage{href-ul}
\hypersetup{
	colorlinks=true,
	linkcolor=blue,
	urlcolor=blue}
\usepackage{geometry}
\geometry{a4paper, top=15mm, left=15mm, right=15mm, bottom=15mm,
	headsep=10mm, footskip=12mm}
%Source: img_0208

% Start the document
\begin{document}
\begin{tikzpicture}(10,3)
	\draw[ultra thick](2,0) --(10,0) -- (10,1.3) --(2,1.3) -- (2,0);
	\draw[fill=black](2,0)-- (0,.6) -- (2,.6) -- (2,0);
	\node at (6,.6) {\large Lösungsblatt von \href{https://www.okuyakl.de}{www.okuyakl.de}};
\end{tikzpicture}
\vspace{0.5 cm}

\noindent{\bf Aufgabe 1. a)}\\
Wir können den Funktionsterm sofort angeben:
$$f^*(x)=0,25x^{-5}$$
\noindent{\bf Aufgabe 1. b)}\\
Definitions- und Wertemenge:
$$\mathbb{D}= \mathbb{R} \setminus \{0\} \qquad  \mathbb{W}= \mathbb{R} \setminus \{0\} $$
Asymptoten sind $x=0$ und $y=0$.
Symmetrie:
$$f^*(-x)=0,25(-x)^{-5} = -0,25x^{-5}=-f^*(x) $$
Die neue Funktion ist punktsymmetrisch.

\noindent{\bf Aufgabe 1. c)}\\
$$
\renewcommand{\arraystretch}{2}
\begin{array}{rccl}
 & x'&=& x+3 \\
\land & y' &=& y+4 \\
\hline
 & x&=& x'-3 \\
\land & y' &=& 0,25x^{-5}+4\\
\hline
\Rightarrow & y' &=& 0,25(x'-3)^{-5}+4\\
& f^{**}(x)&=&  0,25(x-3)^{-5}+4
\end{array}
$$
\begin{minipage}{0.5\textwidth}
\noindent{\bf Aufgabe 2. a)}\\
$$f_2(x)=0,5x^4$$
$$f_3(x)=0,5(x-4)^4+1$$
\noindent{\bf Aufgabe 2. b)}\\
$$\mathbb{D}_{f_3}=\mathbb{R}$$
$$\mathbb{W}_{f_3}=[1;\infty[$$
\end{minipage}
\hfill
\begin{minipage}{0.5\textwidth}
\includegraphics[width=7 cm]{otto208}
\end{minipage}

\noindent{\bf Aufgabe 3. a)}\\
$$\mathbb{D}_{f}=\mathbb{R}\setminus \{1\}$$
$$\mathbb{W}_{f}=[3;\infty[$$
Da die Wertemenge nur positive Zahlen umfasst, schneidet der Graph die $x-$Achse nie.
\vspace{0.5 cm}

\begin{minipage}{0.5\textwidth}
\noindent{\bf Aufgabe 3. b)  c) und d)}\\
$$
\renewcommand{\arraystretch}{2}
\begin{array}{rccl}
 & x'&=& x-3 \\
\land & y' &=& y-4 \\
\hline
 & x&=& x'+3 \\
\land & y' &=& 2(x-1)^{-2}+3-4\\
\hline
\Rightarrow & y' &=& 2(x'+3-1)^{-2}-1\\
& h(x)&=&  2(x+2)^{-2}-1
\end{array}
$$
\end{minipage}
\hfill
\begin{minipage}{0.5\textwidth}
\includegraphics[width=7 cm]{holl208}
\end{minipage}

\noindent{\bf Aufgabe 3. e)}\\
Die waagrechte Asymptote ist bei $y=-1$; die senkrechte ist bei $x=-2$.

\noindent{\bf Aufgabe 3. f)}\\
Gemäß Anleitung zeigen wir:
$$f(-2+a)=2(-2+a+2)^{-2}-1 = 2a^{-2}-1$$
$$f(-2-a)=2(-2-a+2)^{-2}-1 = 2(-a)^{-2}-1 =2a^{-2}-1=f(-2+a) $$
Die Gerade $x=-2$ ist Symmetrieachse.

\noindent{\bf Aufgabe 3. g)}\\
$$
\renewcommand{\arraystretch}{2}
\begin{array}{rcll}
2(x+2)^{-2}-1  &=& 0 &|+1\\
{2 \over (x+2)^2} &=& 1 &|\cdot (x+2)^2\\
2 &=&  (x+2)^2 &|\pm \sqrt{\quad}\\
\pm \sqrt{2} &=& x+2 &|-2\\
\pm\sqrt{2}-2 &=& x \\
x_1 &=&  \sqrt{2}-2 \\
x_2 &=&  -\sqrt{2}-2
\end{array}
$$
Die errechneten Nullstellen stimmen mit denen des Diagramms überein.

\noindent{\bf Aufgabe 4.}\\
a) $$(2+2,3)^0=1 $$
b) $$1+\left({1\over 3}\right)^{-2}=1+3^2 = 10 $$
c) $$1:4^{-2}=1\cdot 4^2=16$$
d) $$2^{x+2}\cdot\left({1\over 2}\right)^{x+2} = 2^{x+2} \cdot 2^{-(x+2)}=2^0=1$$
e) $$\left({ 2 \over a}\right)^3 \cdot (-a)^3=-2^3 =-8$$

\begin{center}
	\includegraphics[width=7 cm]{../../viecher/endcomic.pdf}
	
	Hier geht es zurück zum \href{https://www.okuyakl.de/math/m10pofuL208/aa208.pdf}{Aufgabenblatt}
\end{center}

\end{document}